\begin{trapbox}
	\begin{description}[style=nextline, leftmargin=2.8em, labelsep=0.5em, itemsep=0.35em]
		\item[\textbf{\textcolor{CTrap}{易错点一（旋转轴混淆）}}] 记错公式，将绕 \(x\) 轴旋转的体积公式误用于绕 \(y\) 轴旋转。
		\item[\textbf{\textcolor{CTrap}{正确理解（轴决定截面半径）：}}] 绕哪条轴旋转，截面半径就由垂直该轴的坐标决定。绕 \(x\) 轴时，半径是 \(|y|\)，体积公式含 \(b^{2}\)；若绕 \(y\) 轴，则应使用 \(x\) 的表达式积分，公式结构会互换 \(a,b\) 的角色。
		\item[\textbf{\textcolor{CTrap}{错因分析（死记硬背）：}}] 只记忆体积公式的外形，不思考公式中 \(a,b\) 的几何意义与旋转轴的关系。
	\end{description}
	\begin{center}
		% fig04_axis_compare_fixed.tex
% 用法：在主文件中提前加载 \usepackage{tikz,amsmath}
% 然后 \input{figures/fig04_axis_compare_fixed.tex}
% 说明：这是 TikZ 片段，不包含 \documentclass、\begin{document}、\end{document}。

\begin{center}
	\begin{tikzpicture}[
			x=1cm,
			y=1cm,
			>=stealth,
			line cap=round,
			line join=round,
			every node/.style={font=\small},
			title/.style={font=\bfseries\Large,blue!55!black},
			axis/.style={->,line width=0.8pt,black},
			body/.style={blue!60!black,line width=0.95pt},
			hidden/.style={blue!50!black,dashed,line width=0.65pt},
			dimblue/.style={<->,blue!65!black,line width=0.8pt},
			dimred/.style={<->,red!70!black,line width=0.85pt},
			lab/.style={fill=blue!6,inner xsep=2pt,inner ysep=1pt},
			warn/.style={red!70!black,font=\bfseries,align=center}
		]

		% 手动扩大边界，防止标题、底部说明被裁剪
		\path[use as bounding box] (-3.05,-4.05) rectangle (8.95,2.85);

		% ==================================================
		% 左：绕 x 轴旋转
		% ==================================================
		\begin{scope}[shift={(0,0)}]

			\node[title] at (0,2.35) {绕 $x$ 轴旋转};

			% 旋转椭球体
			\fill[blue!6] (0,0) ellipse (2.20 and 0.95);
			\draw[body] (0,0) ellipse (2.20 and 0.95);

			% 竖向截面虚线
			\draw[hidden] (0,0) ellipse (0.30 and 0.95);

			% x 轴
			\draw[axis] (-2.65,0)--(2.85,0);
			\node at (3.02,0.08) {$x$};

			% 长轴 2a：标签放在线上方，避开下方“三轴半长”
			\draw[dimblue] (-2.20,-1.20)--(2.20,-1.20)
			node[midway,above=3pt,lab] {$2a$};

			% 垂直方向 2b：红色强调，表示截面半径来自 y 方向
			\draw[dimred] (0,-0.95)--(0,0.95)
			node[midway,right=6pt,fill=blue!6,inner xsep=2pt,inner ysep=1pt] {$2b$};

			% 说明文字整体下移，避免和 2a 标注冲突
			\node[align=center] at (0,-2.12)
			{三轴半长：$a,b,b$};

			\node[align=center] at (0,-2.62)
			{$\displaystyle V_x=\frac{4}{3}\pi ab^2$};

		\end{scope}

		% ==================================================
		% 右：绕 y 轴旋转
		% ==================================================
		\begin{scope}[shift={(6.00,0)}]

			% 标题上移，避免与 y 轴标签重叠
			\node[title] at (0,2.35) {绕 $y$ 轴旋转};

			% 旋转椭球体
			\fill[blue!6] (0,0) ellipse (1.05 and 1.65);
			\draw[body] (0,0) ellipse (1.05 and 1.65);

			% 水平截面虚线
			\draw[hidden] (0,0) ellipse (1.05 and 0.25);

			% y 轴：y 标签放到左上侧，避免压标题
			\draw[axis] (0,-1.95)--(0,2.00);
			\node at (-0.20,2.15) {$y$};

			% 水平方向 2a：蓝色表示来自 x 方向
			\draw[dimblue] (-1.05,0)--(1.05,0)
			node[midway,below=5pt,lab] {$2a$};

			% 竖直方向 2b：右移，避免贴住椭球体
			\draw[dimred] (1.48,-1.65)--(1.48,1.65)
			node[midway,right=6pt,fill=blue!6,inner xsep=2pt,inner ysep=1pt] {$2b$};

			% 说明文字下移，和图形保持距离
			\node[align=center] at (0,-2.25)
			{三轴半长：$a,a,b$};

			\node[align=center] at (0,-2.75)
			{$\displaystyle V_y=\frac{4}{3}\pi a^2b$};

		\end{scope}

		% ==================================================
		% 底部提醒
		% ==================================================
		\node[warn] at (3.00,-3.55)
		{旋转轴不同，平方项不同：看清截面半径来自哪个方向};

	\end{tikzpicture}
\end{center}

		\vspace{0.4em}
		{\small 图：椭圆分别绕 \(x\) 轴与 \(y\) 轴旋转的体积公式对比示意图}
	\end{center}
	\medskip
	\begin{description}[style=nextline, leftmargin=2.8em, labelsep=0.5em, itemsep=0.35em]
		\item[\textbf{\textcolor{CTrap}{易错点二（未化标准形式）}}] 遇到非标准方程如 \(4x^{2}+9y^{2}=36\) 时，直接取 \(a=2,b=3\) 代入计算。
		\item[\textbf{\textcolor{CTrap}{正确理解（先归一化）：}}] 必须先将方程化为 \(\frac{x^{2}}{9}+\frac{y^{2}}{4}=1\)，确定正确的半轴 \(a=3,b=2\)。
		\item[\textbf{\textcolor{CTrap}{错因分析（系数处理疏忽）：}}] 误以为方程中 \(x^{2},y^{2}\) 项的系数就是 \(\frac{1}{a^{2}},\frac{1}{b^{2}}\)，忘记等式右边必须为 1。
	\end{description}
	\begin{center}
		% fig05_standardize_fixed.tex
% 用法：在主文件中提前加载 \usepackage{tikz,amsmath}
% 然后 \input{figures/fig05_standardize_fixed.tex}
% 说明：这是 TikZ 片段，不包含 \documentclass、\begin{document}、\end{document}。

\begin{center}
	\begin{tikzpicture}[
			x=1cm,
			y=1cm,
			>=stealth,
			line cap=round,
			line join=round,
			every node/.style={font=\small},
			title/.style={font=\bfseries\Large,blue!55!black},
			bluebox/.style={
					draw=blue!60!black,
					fill=blue!6,
					rounded corners=4pt,
					line width=0.9pt,
					minimum height=1.05cm,
					align=center,
					inner xsep=8pt,
					inner ysep=6pt
				},
			redbox/.style={
					draw=red!70!black,
					fill=red!5,
					rounded corners=4pt,
					line width=0.9pt,
					minimum height=1.05cm,
					align=center,
					inner xsep=8pt,
					inner ysep=6pt
				},
			greenbox/.style={
					draw=green!45!black,
					fill=green!10,
					rounded corners=4pt,
					line width=0.9pt,
					minimum height=1.10cm,
					align=center,
					inner xsep=10pt,
					inner ysep=7pt
				},
			flow/.style={->,red!70!black,line width=1.0pt},
			flowtext/.style={
					font=\small,
					red!70!black,
					fill=white,
					inner xsep=2pt,
					inner ysep=1pt
				},
			warn/.style={red!70!black,font=\bfseries\large,align=center}
		]

		% 防止裁剪
		\path[use as bounding box] (-6.25,-3.10) rectangle (6.25,2.75);

		% 标题
		\node[title] at (0,2.25) {非标准方程：先化标准式，再读半轴};

		% ==================================================
		% 上方三步流程
		% ==================================================
		\node[bluebox,minimum width=3.00cm] (eq1) at (-4.20,0.75)
		{$\displaystyle 4x^2+9y^2=36$};

		\node[bluebox,minimum width=3.25cm] (eq2) at (0,0.75)
		{$\displaystyle \frac{x^2}{9}+\frac{y^2}{4}=1$};

		\node[redbox,minimum width=2.95cm] (ans) at (4.15,0.75)
		{$\displaystyle a=3,\ b=2$};

		% 箭头与说明
		\draw[flow] (eq1.east) -- (eq2.west)
		node[midway,above=7pt,flowtext] {同除以 $36$};

		\draw[flow] (eq2.east) -- (ans.west)
		node[midway,above=7pt,flowtext] {看分母};

		% ==================================================
		% 下方：错误读法 与 正确代入
		% ==================================================
		\node[redbox,minimum width=4.25cm] (bad) at (-2.65,-1.05)
		{错误读法：$\displaystyle a=2,\ b=3$};

		% 红叉
		\draw[red!70!black,line width=1.25pt] (-4.10,-1.42)--(-1.20,-0.68);
		\draw[red!70!black,line width=1.25pt] (-4.10,-0.68)--(-1.20,-1.42);

		\node[greenbox,minimum width=4.35cm] (good) at (2.70,-1.05)
		{$\displaystyle V=\frac{4}{3}\pi\cdot 3\cdot 2^2=16\pi$};

		% ==================================================
		% 底部提醒
		% ==================================================
		\node[warn] at (0,-2.35)
		{等式右边必须先变成 $1$，不能直接从系数读半轴};

	\end{tikzpicture}
\end{center}

		\vspace{0.4em}
		{\small 图：非标准方程先化标准式流程图}
	\end{center}
	\medskip
	\begin{description}[style=nextline, leftmargin=2.8em, labelsep=0.5em, itemsep=0.35em]
		\item[\textbf{\textcolor{CTrap}{易错点三（半轴与轴长混淆）}}] 将轴长（如长轴长 10）直接作为 \(a=10\) 代入公式。
		\item[\textbf{\textcolor{CTrap}{正确理解（半轴长定义）：}}] 椭圆标准方程中的 \(a\) 是半轴长，即轴长的一半。长轴长为 \(2a\)，故 \(a=5\)。
		\item[\textbf{\textcolor{CTrap}{错因分析（概念模糊）：}}] 对椭圆几何参数定义记忆不清，混淆了“轴”与“半轴”。
	\end{description}
	\begin{center}
		% fig06_semiaxis_warning_fixed.tex
% 用法：在主文件中提前加载 \usepackage{tikz,amsmath}
% 然后 \input{figures/fig06_semiaxis_warning_fixed.tex}
% 说明：这是 TikZ 片段，不包含 \documentclass、\begin{document}、\end{document}。

\begin{center}
	\begin{tikzpicture}[
			x=1cm,
			y=1cm,
			>=stealth,
			line cap=round,
			line join=round,
			every node/.style={font=\small},
			title/.style={font=\bfseries\Large,blue!55!black},
			axis/.style={->,line width=0.8pt,black},
			body/.style={blue!60!black,line width=0.95pt},
			dim/.style={<->,blue!65!black,line width=0.85pt},
			semiaxis/.style={->,red!70!black,line width=1.05pt},
			labBlue/.style={
					fill=blue!6,
					inner xsep=3pt,
					inner ysep=1.5pt
				},
			labRed/.style={
					fill=blue!6,
					text=red!70!black,
					inner xsep=3pt,
					inner ysep=1.5pt
				},
			warn/.style={red!70!black,font=\bfseries\large,align=center}
		]

		% ==================================================
		% 手动扩大边界，防止裁剪
		% ==================================================
		\path[use as bounding box] (-4.65,-3.85) rectangle (4.65,2.95);

		% ==================================================
		% 标题
		% ==================================================
		\node[title] at (0,2.48) {公式里用的是半轴长，不是轴长};

		% ==================================================
		% 椭圆本体
		% ==================================================
		\fill[blue!6] (0,0) ellipse (2.75 and 1.15);
		\draw[body] (0,0) ellipse (2.75 and 1.15);

		% ==================================================
		% 坐标轴
		% ==================================================
		\draw[axis] (-3.25,0)--(3.45,0);
		\node at (3.62,0.08) {$x$};

		\draw[axis] (0,-1.62)--(0,1.65);
		\node at (0.18,1.78) {$y$};

		% ==================================================
		% 蓝色：轴长
		% ==================================================
		% 水平方向轴长 2a
		\draw[dim] (-2.75,-1.42)--(2.75,-1.42);
		\node[labBlue] at (0,-1.72) {轴长 $2a$};

		% 竖直方向轴长 2b
		\draw[dim] (-3.05,-1.15)--(-3.05,1.15);
		\node[labBlue] at (-3.78,0) {轴长 $2b$};

		% ==================================================
		% 红色：半轴
		% ==================================================
		% 半轴 a：标签放到右侧，不和半轴 b 挤在一起
		\draw[semiaxis] (0,0)--(2.75,0);
		\node[labRed] at (1.92,0.34) {半轴 $a$};

		% 半轴 b：标签放到左上侧，和半轴 a 明确分开
		\draw[semiaxis] (0,0)--(0,1.15);
		\node[labRed] at (-0.55,0.78) {半轴 $b$};

		% ==================================================
		% 底部提醒文字
		% ==================================================
		\node[warn] at (0,-2.38)
		{例如长轴长为 $10$，代入公式的是 $a=5$};

		\node[warn] at (0,-2.88)
		{不是 $a=10$};

		\node[align=center] at (0,-3.42)
		{$\displaystyle V=\frac{4}{3}\pi ab^2$ 中的 $a,b$ 都是从中心到端点的长度};

	\end{tikzpicture}
\end{center}

		\vspace{0.4em}
		{\small 图：半轴长不是轴长示意图}
	\end{center}
	\medskip
	\textbf{\textcolor{CTrap}{新增易错点四（沿 y 方向切片误认为绕 y 轴旋转）}}
	\begin{description}[style=nextline, leftmargin=2.8em, labelsep=0.5em, itemsep=0.35em]
		\item[\textbf{错法：}] 看到沿 \(y\) 方向积分 \(\int_{-b}^{b}S(y)\,dy\)，便误认为旋转轴变成了 \(y\) 轴，或误将截面当作圆。
		\item[\textbf{正确：}] 旋转轴从始至终都是 \(x\) 轴。\textbf{沿 \(y\) 切片}是固定 \(y=c\)，得到 \(xz\) 平面内的椭圆，面积为 \(S(c)=\pi ab(1-c^{2}/b^{2})\)；\textbf{绕 \(y\) 轴旋转}则会生成另一个椭球，体积一般为 \(\frac43\pi a^{2}b\)。两件事不能混淆。
		\item[\textbf{边界提醒：}] 在 \(y=\pm b\) 的两个端点，椭圆截面退化成一个点，不应说所有高度的截面都是非退化椭圆。
	\end{description}

	\medskip
	\textbf{\textcolor{CTrap}{新增易错点五（球体伸缩倍数颠倒）}}
	\begin{description}[style=nextline, leftmargin=2.8em, labelsep=0.5em, itemsep=0.35em]
		\item[\textbf{错法：}] 半径 \(b\) 的球变成 \(x\) 半轴为 \(a\) 的椭球时，误乘 \(b/a\)，或者把三个方向都乘 \(a/b\)。
		\item[\textbf{正确：}] 仅 \(x\) 方向从 \(b\) 变为 \(a\)，所以体积乘 \(a/b\)；\(y,z\) 方向保持 \(b\) 不变。坐标替换为 \(x=(a/b)u\)，反代时才是 \(u=(b/a)x\)。
	\end{description}

\end{trapbox}
